\documentclass[11pt]{article}
\usepackage[margin=1in]{geometry}
\usepackage{amsmath,amssymb,amsthm,mathtools,microtype}
\usepackage[hidelinks]{hyperref}
\newtheorem{theorem}{Theorem}
\newtheorem{lemma}[theorem]{Lemma}
\newtheorem{remark}[theorem]{Remark}
\newcommand{\R}{\mathbb R}
\newcommand{\Tr}{\operatorname{Tr}}
\newcommand{\dist}{\operatorname{dist}}
\title{The one-dimensional endpoint Hardy--Lieb--Thirring inequality on arbitrary domains}
\author{Autonomous proof attempt for arXiv:1102.4394}
\date{11 August 2026}

\begin{document}
\maketitle

\begin{abstract}
Frank and Loss proved a Hardy--Lieb--Thirring inequality on arbitrary domains
in dimension one for moments $\gamma>1/2$ and wrote that the endpoint
$\gamma=1/2$ should also hold.  We prove that endpoint.  A finite interval is
split at its midpoint.  The two Dirichlet halves are controlled by the known
critical half-line endpoint inequality.  Removing the midpoint condition is a
codimension-one form extension, so interlacing leaves only the lowest
eigenvalue.  The source paper's pointwise Hardy--Sobolev inequality controls
that eigenvalue.  Summing over interval components gives a universal constant,
explicitly $C_{1/2,0}+2\leq3.185$ in terms of the published half-line constant.
\end{abstract}

\section{Statement}

Let $\Omega\subsetneq\R$ be open and put
$$d_\Omega(x)=\dist(x,\Omega^c).$$
In dimension one the averaged distance $D_\Omega$ in the source paper equals
$d_\Omega$.  Let $H_{\Omega,V}$ be the Friedrichs operator associated with
the quadratic form
\begin{equation}\label{eq:form}
 \mathfrak q_{\Omega,V}[u]
 =\int_\Omega\left(|u'|^2-\frac{|u|^2}{4d_\Omega^2}
                         +V|u|^2\right)dx.
\end{equation}
The source proves
$\Tr(H_{\Omega,V})_-^\gamma\le L_{1,\gamma}
\int_\Omega V_-^{\gamma+1/2}$ for $\gamma>1/2$ and asks immediately
after Theorem~1.5 whether it remains true at $\gamma=1/2$.

\begin{theorem}\label{thm:main}
There is a universal constant $L_{1,1/2}$ such that, for every open
$\Omega\subsetneq\R$ and every real $V\in L^1(\Omega)$,
\begin{equation}\label{eq:main}
 \Tr(H_{\Omega,V})_-^{1/2}
 \le L_{1,1/2}\int_\Omega V_-(x)\,dx.
\end{equation}
One may take
$$L_{1,1/2}=C_{1/2,0}+2\leq3.185,$$
where $C_{1/2,0}\leq1.185$ is the upper bound in the critical half-line
inequality of Ekholm and Frank.
\end{theorem}

\section{Two inputs}

The first input is Ekholm--Frank, Theorem~2.1 in
arXiv:math/0611247.  In the special case $\alpha=0$, $\gamma=1/2$ it says
that, for $W\ge0$ on $\R_+=(0,\infty)$,
\begin{equation}\label{eq:halfline}
 \Tr\left(-\frac{d^2}{dr^2}-\frac1{4r^2}-W\right)_-^{1/2}
 \le C_{1/2,0}\int_0^\infty W(r)\,dr,
\end{equation}
with the Dirichlet condition at the singular endpoint.

The second input is Corollary~2.2 of the source paper at $q=2$.  If $I$ is
an interval and
\begin{equation}\label{eq:hardyform}
 \mathfrak h_I[u]=\int_I\left(|u'|^2-\frac{|u|^2}{4d_I^2}\right)dx,
 \qquad d_I(x)=\dist(x,I^c),
\end{equation}
then
\begin{equation}\label{eq:pointwise}
 \|u\|_\infty^4\le16\,\mathfrak h_I[u]\,\|u\|_2^2.
\end{equation}
The source states the bound first on $C_c^\infty(I)$; it extends to the
closed form domain by density.  In particular, evaluation at an interior
point is continuous in the form norm.

\section{A finite interval}

Fix $I=(a,b)$, let $m=(a+b)/2$, and write $W=V_-\ge0$ and
$A_I=\int_IW$.  Since positive parts of $V$ only raise the operator, it
suffices to study
$$H_I=-\frac{d^2}{dx^2}-\frac1{4d_I(x)^2}-W.$$
Let $H_I^D$ be the form restriction obtained by imposing $u(m)=0$.  It is the
orthogonal direct sum of the operators on $(a,m)$ and $(m,b)$ with a
Dirichlet condition at $m$.  On the left half $d_I(x)=x-a$, and on the right
half $d_I(x)=b-x$.  After translation or reflection, each half is therefore
the critical half-line form on $(0,R)$, $R=(b-a)/2$, with an extra Dirichlet
condition at $R$.

Zero extension at $R$ embeds this finite-interval Dirichlet form domain into
the half-line form domain.  Extend the corresponding half-potential by zero
beyond $R$.  The min--max principle and \eqref{eq:halfline}, applied to both
halves, give
\begin{equation}\label{eq:split}
 \Tr(H_I^D)_-^{1/2}\le C_{1/2,0}A_I.
\end{equation}

The condition $u(m)=0$ has codimension one.  If
$\mu_1\le\mu_2\le\cdots$ and $\nu_1\le\nu_2\le\cdots$ are respectively
the negative eigenvalues of $H_I$ and $H_I^D$, with the lists continued by
the bottom of the nonnegative spectrum when needed, form interlacing gives
\begin{equation}\label{eq:interlace}
 \mu_j\le\nu_j\le\mu_{j+1}.
\end{equation}
Consequently
\begin{equation}\label{eq:rankone}
 \Tr(H_I)_-^{1/2}
 \le(-\mu_1)_+^{1/2}+\Tr(H_I^D)_-^{1/2}.
\end{equation}

It remains to estimate the single possible extra eigenvalue.  For
$\|u\|_2=1$, \eqref{eq:pointwise} implies
\begin{align*}
 \mathfrak h_I[u]-\int_IW|u|^2
 &\ge \mathfrak h_I[u]-A_I\|u\|_\infty^2\\
 &\ge \mathfrak h_I[u]-4A_I\mathfrak h_I[u]^{1/2}
 \ge-4A_I^2.
\end{align*}
Thus $\mu_1\ge-4A_I^2$, and \eqref{eq:split}--\eqref{eq:rankone} yield
\begin{equation}\label{eq:component}
 \Tr(H_I)_-^{1/2}\le(C_{1/2,0}+2)A_I.
\end{equation}

\section{Arbitrary open sets}

Every open proper subset of $\R$ is a countable disjoint union of finite
intervals and rays.  On a ray the distance to the complement is exactly the
distance to its finite endpoint, so \eqref{eq:halfline} applies directly and
gives the stronger constant $C_{1/2,0}$.  On every finite component use
\eqref{eq:component}.  The operator on $\Omega$ is the orthogonal direct sum
of the component operators.  Summing their eigenvalue moments gives
\begin{align*}
 \Tr(H_{\Omega,-V_-})_-^{1/2}
 &\le(C_{1/2,0}+2)\sum_I\int_I V_-\\
 &=(C_{1/2,0}+2)\int_\Omega V_-.
\end{align*}
This proves \eqref{eq:main} for bounded compactly supported $V_-$.  Monotone
truncation and standard form convergence give the result for every
$V_-\in L^1(\Omega)$, and monotonicity restores the positive part of $V$.
\qed

\section*{Scope and source note}

This settles the explicit $N=1$, $\gamma=1/2$ sentence on source PDF page 5.
It does not settle the paper's separate question about a $p$-Hardy--Sobolev
inequality for $1<p<2$.  The proof uses only the source's Corollary~2.2 and
Tomas Ekholm and Rupert L. Frank, \emph{Lieb--Thirring inequalities on the
half-line with critical exponent}, arXiv:math/0611247, Theorem~2.1.

\end{document}
